\section{问题二：变物性条件下的热湿耦合}
\subsection{统一物性与耦合方程}
问题二建立整个烘干过程的变物性模型，全程统一使用附录3经验式，本问给出前3 h的结果。因此保留问题一的固定几何、原始初值与边界结构，从$t=0$重新积分。局部物性为
\begin{equation}\label{eq_prop3}
\begin{aligned}
\rho(C)&=650+128C,\\
c_p(C)&=1450+2736\frac{C}{1+C},\\
k(C)&=0.21+0.38\frac{C}{1+C},\\
D(C,T)&=2.4\times10^{-3}\exp\!\left(-\frac{0.45}{C}\right)
\exp\!\left[-\frac{3850}{T+273.15}\right].
\end{aligned}
\end{equation}
其中$\rho$、$c_p$、$k$和$D$的单位分别为kg/m$^3$、J/(kg$\cdot$K)、W/(m$\cdot$K)和m$^2$/s，$C$为干基含水率，单位为kg/kg。温度场变量$T$按$^\circ$C表示，扩散系数指数中的$T+273.15$为药材局部绝对温度，单位为K。

问题二仍采用与问题一相同的固定圆柱环层作为控制体，环层几何和水分守恒结构不变。附录3将热物性与扩散系数改为局部温湿状态的函数，$\rho(C)c_p(C)$仅作为有效体积热容量$B(C)$，固定域水分守恒仍按参考干物质骨架计量。根据假设（3），仍不显式引入汽化潜热和水分携焓，故在相同守恒结构下更新本构关系并联合求解。将式\eqref{eq_prop3}代入该环层模型，得到
\begin{equation}\label{eq_coupled}
\begin{aligned}
B(C)T_t&=\frac1r\partial_r[rk(C)T_r],\\
C_t&=\frac1r\partial_r[rD(C,T)C_r],\\
B(C)&=\rho(C)c_p(C).
\end{aligned}
\end{equation}
初始条件沿用式\eqref{eq_initial}，中心仍采用零梯度条件；表面边界沿用式\eqref{eq_robin}并将$k,D$替换为当前表面物性。题目附录3未另给交换系数，按假设（4）继续使用前述$h,h_m$。

式\eqref{eq_coupled}中的耦合有两个方向：$C$改变$B$和$k$，从而改变温度传播；$T$改变$D$，从而改变水分迁移。进一步展开水分扩散项可得
\begin{equation}\label{eq_moisture_expanded_q2}
C_t=D\left(C_{rr}+\frac{C_r}{r}\right)+D_C(C_r)^2+D_TT_rC_r.
\end{equation}
式\eqref{eq_moisture_expanded_q2}中的$D_TT_rC_r$由空间变化的扩散系数产生，不是另加的热扩散经验项。采用原扩散通量形式离散，可同时保留上述非线性项。

\subsection{有效热容量与干物质守恒核算}
问题二沿用固定参考干物质骨架，水分方程中的参考干密度保持恒定。经验$\rho(C)$仅用于构造$B(C)$，不能在水分守恒中又取$s=\rho(C)/(1+C)$并随时间更新，否则会与固定体积中干物质量不变的假设冲突。

同样，$B(C)T_t$不能直接替换为$\partial_t[B(C)T]$。取初始参考温度$T_{\mathrm{ini}}=28\ ^\circ$C，有
\begin{equation}\label{eq_B_identity}
\partial_t[B(C)(T-T_{\mathrm{ini}})]=B(C)T_t+B'(C)(T-T_{\mathrm{ini}})C_t.
\end{equation}
式\eqref{eq_B_identity}中的第二项反映所定义状态量随组成变化而改变。本文采用以温度形式表达的有效热方程，在收支检验中保留修正项，但这并不代表已经描述完整多相系统的能量交换。

\subsection{联合求解与结果}
沿用问题一的环形控制体，热导率与扩散系数均取相邻节点的调和平均。将$\boldsymbol y=(T_0,\ldots,T_N,C_0,\ldots,C_N)^{\mathsf T}$作为联合状态，其中$T_i$和$C_i$分别表示第$i$个节点的温度和干基含水率；这里的下标$0$表示节点编号，不表示参考温度。在每次右端函数求值及隐式迭代中，按局部温湿状态更新物性。每个节点只与相邻节点的两个场相联系，利用这一稀疏结构进行BDF积分。生产网格采用$G_4$，相对容差$10^{-9}$，温度、含水率绝对容差分别为$10^{-8}$、$10^{-10}$，最大步长1 s。

表\ref{tab_q2T}和表\ref{tab_q2C}给出每0.5 h的论文摘录结果，图\ref{fig_q2}展示相应径向剖面，用于比较烘干过程中不同位置的升温与失水情况；完整的逐秒、逐0.1 cm结果另存于\texttt{result2.xlsx}中。
\begin{table}[htbp]\centering\small
\caption{3小时内药材的温度（$^\circ$C）}\label{tab_q2T}
\renewcommand{\arraystretch}{1.12}
\begin{tabular}{cccccc}\toprule
 & \multicolumn{5}{c}{到药材中心的距离/cm}\\
时间/h & 0 & 0.5 & 1.0 & 1.5 & 2.0\\\midrule
0.5 & 32.1896 & 32.3824 & 32.9660 & 33.9608 & 35.4132\\
1 & 40.3822 & 40.5542 & 41.0609 & 41.8783 & 43.0000\\
1.5 & 45.8471 & 45.9351 & 46.1932 & 46.6055 & 47.1397\\
2 & 48.4503 & 48.4881 & 48.5985 & 48.7736 & 49.0033\\
2.5 & 49.4670 & 49.4793 & 49.5136 & 49.5645 & 49.6617\\
3 & 49.8494 & 49.8553 & 49.8746 & 49.9096 & 49.9645\\
\bottomrule\end{tabular}
\end{table}

\begin{table}[htbp]\centering\small
\caption{3小时内药材的水分浓度（kg/kg）}\label{tab_q2C}
\renewcommand{\arraystretch}{1.12}
\begin{tabular}{cccccc}\toprule
 & \multicolumn{5}{c}{到药材中心的距离/cm}\\
时间/h & 0 & 0.5 & 1.0 & 1.5 & 2.0\\\midrule
0.5 & 2.5499 & 2.5489 & 2.5256 & 2.3257 & 1.6486\\
1 & 2.5257 & 2.4948 & 2.3578 & 2.0230 & 1.4711\\
1.5 & 2.3861 & 2.3256 & 2.1344 & 1.8020 & 1.3476\\
2 & 2.1708 & 2.1084 & 1.9236 & 1.6259 & 1.2311\\
2.5 & 1.9566 & 1.9006 & 1.7360 & 1.4720 & 1.1166\\
3 & 1.7662 & 1.7165 & 1.5702 & 1.3333 & 1.0081\\
\bottomrule\end{tabular}
\end{table}

\figpair{q2_T_profiles.pdf}{q2_C_profiles.pdf}{前3 h变物性模型的温度与含水率剖面。温度逐渐接近烘房水平，水分仍保持中心高、表面低的径向差异。}{fig_q2}

3 h时中心温度为49.8494 $^\circ$C，已接近烘房温度，但中心含水率仍为1.7662 kg/kg，远高于0.15 kg/kg的干燥要求。由此可见，温度接近稳定不能作为停止烘干的条件。

问题一和问题二在相同时间的结果差异主要源于物性设定和耦合方式不同：问题一将温度场与干基含水率场分别求解，问题二则对两个场进行双向耦合求解。问题二不能在1800 s处读取问题一末态后才切换参数，否则前0.5 h，即$0\sim1800$ s的过程不再满足题目要求的统一经验式。按题目输出要求，从1 s到10800 s以1 s、0.1 cm采样，温度场和干基含水率场各含10800$\times$21个数值，合计453600个数值。

\subsection{升温与失水对扩散系数的竞争}
由式\eqref{eq_prop3}直接求导可得
\begin{equation}\label{eq_logD}
\frac{\partial\ln D}{\partial C}=\frac{0.45}{C^2}>0,\qquad
\frac{\partial\ln D}{\partial T}=\frac{3850}{(T+273.15)^2}>0.
\end{equation}
升温倾向于增大扩散系数，失水则使其减小，两种作用沿实际轨迹竞争。固定$C=2.55$时，由28升至50 $^\circ$C使$D$增至约2.388倍；固定$T=50\ ^\circ$C时，$C$从2.55降至0.15使$D$降至$C=2.55,T=50\ ^\circ$C时$D$值的约5.94\%。这些是对经验式的条件性代入，不是独立温度控制实验。

数值结果中前3 h局部$D$约在$5.5476\times10^{-9}$至$1.2522\times10^{-8}$ m$^2$/s之间变化。因各位置升温和失水程度不同，不能用单一常数代表整个时空场。与固定为初始值$D_0$的对照算例相比，3 h内干基含水率最大差为0.3981 kg/kg，说明将$D$简化为常数会明显改变干基含水率场，正式计算应保留变化的$D$。至于其对烘干时长的影响，还需要在问题三中通过临界时间进一步验证。
\FloatBarrier
